\section{Conclusões}

\begin{itemize}
    \item A base analisada tem 1.647 gravações de 22 participantes, todas com formato válido, coincidindo com o artigo; os identificadores s18 e s20 não existem e algumas combinações têm 14 ou 16 gravações.
    \item No par 20, o nível da amplitude é dominado pelo participante (e possivelmente pela sessão de coleta), e não pela atividade. As atividades de transição (SD, GU, LD, FA) tendem a ter variação concentrada em um trecho, enquanto WA tem variação distribuída por toda a gravação.
    \item As cinco características se sobrepõem entre atividades e, isoladamente, não separam as cinco classes; o pico a pico é muito sensível a valores extremos e a média das diferenças absolutas diferencia mais os participantes.
    \item Os mapas de calor mostram pares sem informação útil (par 1 constante, par 2 constante, pares 28-38 próximos de zero), coerentes com as subportadoras nulas e piloto removidas pelos autores.
    \item A recombinação de IMFs por EEMD preserva bem o nível e a dispersão dos sinais reais, mas aumenta a variação entre registros consecutivos e pode gerar valores negativos e fora da faixa real. A configuração da EEMD precisou ser ajustada (separate\_trends = True) para evitar que o nível do sinal fosse somado duas vezes - uma fragilidade também presente no código original.
\end{itemize}
